
\subsubsection{3.1 Spectral Memory Horizon}

State Space Models maintain a hidden state h\_t that evolves according to discretized dynamics:

h\_t = A\_bar \textit{ $h_{t-1}$ + B\_bar } x\_t
y\_t = C \textit{ h\_t + D } x\_t

where A\_bar = exp(Delta \textit{ A) is the discretized transition matrix. For diagonal A (as in Mamba), eigenvalues are given by lambda\_i = exp(Delta } a\_i) where a\_i = -exp(A\_log\_i).

\textbf{Definition (Spectral Memory Horizon):}

H\_spec = -1 / log|lambda\_max|

where lambda\_max = max\_i |lambda\_i| is the eigenvalue with largest magnitude across all SSM layers. This quantity represents the theoretical number of timesteps for which information can persist in the slowest-decaying eigenmode before falling to 1/e of its original magnitude.

For Mamba-1.4B with 48 layers, H\_spec is computed directly from the pretrained A\_log parameters.

\subsubsection{3.2 Eigenvalue Preservation Measurement}

Projection-only LoRA is applied to in\_proj and x\_proj matrices only. Eigenvalue preservation is measured through:

\begin{enumerate}
\item \textbf{Relative H\_spec Change:} |delta H\_spec| = |(H\_spec\_post - H\_spec\_pre) / H\_spec\_pre| x 100\%
\item \textbf{Eigenvalue Correlation:} Pearson correlation between pre- and post-training eigenvalue vectors
\item \textbf{A\_log Maximum Difference:} max|A\_log\_post - A\_log\_pre| across all layers
\end{enumerate}

\subsubsection{3.3 Eigenmode Energy Redistribution}

The slow mode fraction is defined as:

$f_\text{slow} = \sum_{i: |\lambda_i| > 0.99} |h_t[i]|^2 / \sum_i |h_t[i]|^2$

Energy redistribution is measured as:

delta E = |f\_slow\_post - f\_slow\_pre|

A threshold of delta E > 0.1 nats indicates meaningful energy redistribution.

\subsubsection{3.4 Experimental Design}

Four sequential sub-hypotheses were tested:

\begin{table*}[htbp]
\centering
\resizebox{\dimexpr 2\columnwidth+\columnsep\relax}{!}{%
\begin{tabular}{lllllll}
\toprule
Hypothesis & Type & Gate & Metric & Threshold &  &  \\
\midrule
H-E1 & Existence & MUST\_WORK & CV(H\_spec) & < 0.3 &  &  \\
H-M1 & Mechanism & MUST\_WORK & Degradation ratio & > 1.1 &  &  \\
H-M2 & Mechanism & MUST\_WORK &  & delta H\_spec &  & < 10\% \\
H-M3 & Mechanism & SHOULD\_WORK & delta E & > 0.1 &  &  \\
\bottomrule
\end{tabular}
}
\end{table*}

H-M4 (discriminative MQAR test) was planned but not executed.
